\documentclass[11pt]{article}

\usepackage[utf8]{inputenc}
\usepackage[T1]{fontenc}
\usepackage{lmodern}
\usepackage{geometry}
\usepackage{amsmath,amssymb,amsfonts,amsthm,mathtools}
\usepackage{bm}
\usepackage{enumitem}
\usepackage{microtype}
\usepackage{hyperref}
\usepackage{titlesec}
\usepackage{booktabs}
\usepackage{array}
\usepackage{setspace}
\usepackage{mathrsfs}
\usepackage{physics}

\geometry{margin=1in}
\hypersetup{colorlinks=true, linkcolor=black, urlcolor=blue, citecolor=black}
\setlist[enumerate]{topsep=0.4em,itemsep=0.35em}
\setlist[itemize]{topsep=0.3em,itemsep=0.25em}
\titleformat{\section}{\large\bfseries}{\thesection.}{0.5em}{}
\titleformat{\subsection}{\normalsize\bfseries}{\thesubsection.}{0.5em}{}
\setstretch{1.06}

\newcommand{\Aether}{\AE{}ther}
\newcommand{\Aflow}{\AE{}ther-flow}
\newcommand{\TheoryName}{\Aflow{} Interpretation of Relativity}
\newcommand{\TheoryProgram}{\Aether{} / \Aflow{} framework}

\newcommand{\CompletedMetric}{g^{\sharp}}
\newcommand{\MatterSpace}{\Psi}

\newcommand{\Car}{\mathrm{car}}
\newcommand{\Cur}{\mathrm{cur}}
\newcommand{\Hom}{\mathrm{hom}}

\newcommand{\ForSpace}[1]{\mathcal I_{#1}^{\mathrm{for}}}
\newcommand{\PrimShell}{\mathcal S_{\mathrm{prim}}}
\newcommand{\Reduction}[1]{\mathcal R_{#1}}
\newcommand{\CurrentCarrier}[1]{\mathfrak C_{#1}^{\mathrm{cur}}}
\newcommand{\EqCarrier}[1]{\mathfrak C_{#1}^{\mathrm{eq}}}
\newcommand{\MetricOut}{\mathscr G_{\mathrm{out}}}
\newcommand{\MatterOut}{\mathscr T_{\mathrm{out}}}

\newcommand{\OriginPackage}{\mathfrak O^{\mathrm{sub,orig}}}
\newcommand{\GroundPackage}{\mathfrak G^{\mathrm{sub,grd}}}
\newcommand{\GroundCore}{\mathfrak G^{\mathrm{sub,grd,core}}}
\newcommand{\PrecursorPackage}{\mathfrak P^{\mathrm{sub,prec}}}
\newcommand{\LiftPackage}{\mathfrak L^{\mathrm{sub,lift}}}
\newcommand{\ShellRestriction}{\mathfrak S^{\mathrm{sub,sh}}}

\newtheorem{definition}{Definition}
\newtheorem{proposition}{Proposition}
\newtheorem{remark}{Remark}
\newtheorem{corollary}{Corollary}
\newtheorem{theorem}{Theorem}

\title{\textbf{The \TheoryName{}}\\[0.4em]
\large Primitive-Reservoir Observer-Reduction\\
\large Precursor-Generating Substrate Origin Dynamics\\
\large Necessity / Uniqueness from\\
\large Origin-Generating Substrate Ground Dynamics}
\author{}
\date{}

\begin{document}

\maketitle

\begin{abstract}
The current benchmark-facing frontier on the active primitive-reservoir observer-reduction line is no longer the lift-generating substrate precursor package \(\PrecursorPackage\) itself. The precursor-forcing theorem has already spent that move: once precursor-generating substrate origin dynamics are fixed, the precursor package is canonical and unique. The next honest benchmark-facing target is therefore the origin package \(\OriginPackage\) that supports precursor scope.

The recorded origin-from-ground theorem already showed that \(\OriginPackage\) can be produced canonically from a deeper origin-generating substrate ground class \(\GroundPackage\). The present note sharpens that existence statement to the forcing verdict required by the routing layer. For each sector it isolates the benchmark-relevant fiber-minimizing exact-shell ground core: the ground-to-origin projection and the recorded origin-, precursor-, lift-, and substrate-fiber minimizer images, the exact ground shell and its observer-data and shell-response readouts, the shell chart to the same reservoir body, the support shell and zero core, the zero/hidden split in the same equation variables, the hidden charge and exact zero-response realization, and the fixed-data solvability and coercive control carried on that core. The theorem proves that every ground-faithful presentation inducing the fixed origin package determines exactly the same ground core up to canonical ground-faithful equivalence.

In that precise sense the precursor-generating substrate origin dynamics are no longer merely available inside a deeper class. They are necessary and unique at benchmark-facing scope as the projected exact-shell output of origin-generating substrate ground dynamics. The benchmark boundary remains fixed throughout: the one operative metric is still exactly \(\CompletedMetric\), the ordinary matter route is unchanged, there is no second operative metric, and the low-energy matter law is not enlarged. The honest burden therefore moves below origin scope to the ground layer itself.
\end{abstract}

\section{Introduction}

The active derivational program of \TheoryName{} has now reached the point where another precursor relay would be benchmark neutral. The current precursor-forcing theorem already removed the lift-generating substrate precursor package \(\PrecursorPackage\) as the live stopping point on the active primitive-reservoir observer-reduction line. Once that theorem is on record, the next honest benchmark-facing move must go one level lower and address the precursor-generating substrate origin dynamics that support precursor scope.

There are two disciplined ways to do that. One can derive the origin package from still more primitive explicit substrate structure, or one can prove that the origin package is already forced inside a sharper deeper class. The recorded origin-from-ground theorem supplied the first ingredient needed for the second route: it exhibited the entire origin package as a projected exact-shell transport of a deeper ground class. What it did not yet state sharply enough for the current routing burden is the corresponding necessity / uniqueness verdict at origin scope.

The present note records exactly that stronger claim. It does not change the benchmark exact-closure package. It does not introduce a second operative metric or a new low-energy matter law. It only proves that, once the deeper ground package is fixed, no independent benchmark-facing freedom remains at precursor-generating origin scope.

\paragraph{Framework and claim boundary.}
\TheoryProgram{} denotes the broader \Aether{} / \Aflow{} framework built on the claims that the \Aether{} is the underlying four-dimensional substrate of reality, the \Aflow{} is its intrinsic ordered motion, observed three-dimensional space is the local experiential slice of that deeper substrate, S-time is the experienced order of change arising from matter, light, and the \Aflow{}, and observed expansion is the three-dimensional appearance of deeper four-dimensional ordered motion. Within that framework, \TheoryName{} denotes the exact-closure theory in which the effective gravitational dynamics are adopted to be exactly Einsteinian with universal matter coupling. In the active sequence, \emph{adoption} means use of that established relativistic dynamics without claiming substrate derivation, while \emph{derivation} is reserved for a first-principles recovery from explicit substrate variables.

The present note stays strictly inside that boundary. It does not claim that final substrate microphysics has already been identified. Its narrower theorem is only that, on the active primitive-reservoir observer-reduction line, the precursor-generating substrate origin dynamics are themselves a forced and unique output of the recorded ground class that sits immediately below them.

\section{Fixed Primitive Continuation Data}

\begin{definition}[Recorded primitive continuation data]
\label{def:fixed}
Fix the following continuation data from the active primitive-reservoir observer-reduction line.
\begin{enumerate}
    \item For each sector label \(\sigma\in\{\Car,\Cur,\Hom\}\), a primitive forcing shell
    \[
    \ForSpace{\sigma},
    \]
    a visible reduction map
    \[
    \Reduction{\sigma}:\ForSpace{\sigma}\to X_{\sigma},
    \]
    and shell-level current and equilibrium carriers
    \[
    \CurrentCarrier{\sigma}:\ForSpace{\sigma}\to Z_{\sigma}^{\mathrm{cur}},
    \qquad
    \EqCarrier{\sigma}:\ForSpace{\sigma}\to Z_{\sigma}^{\mathrm{eq}}.
    \]
    \item The product shell
    \[
    \PrimShell
    \equiv
    \ForSpace{\Car}\times \ForSpace{\Cur}\times \ForSpace{\Hom}.
    \]
    \item The recorded metric and ordinary-matter outputs
    \[
    \MetricOut:\PrimShell\to\{\text{observer metrics}\},
    \qquad
    \MatterOut:\PrimShell\times\MatterSpace\to\{\text{observer stress tensors}\},
    \]
    satisfying
    \begin{equation}
    \MetricOut(\mathbf w)=\CompletedMetric,
    \qquad
    \MatterOut(\mathbf w,\psi)
    =
    T_{\mu\nu}^{\mathrm{matter}}[\CompletedMetric,\psi]
    \label{eq:fixedoutputs}
    \end{equation}
    for every \(\mathbf w\in\PrimShell\) and every matter configuration \(\psi\in\MatterSpace\).
\end{enumerate}
\end{definition}

\begin{remark}
Definition \ref{def:fixed} is not reopened here. The primitive continuation data, the one operative metric, and the ordinary matter route remain fixed. The present note only addresses the remaining benchmark-facing freedom at origin scope.
\end{remark}

\section{Recorded Origin Target}

\begin{definition}[Recorded precursor-generating substrate origin package]
\label{def:orig}
Fix, for each sector \(\sigma\in\{\Car,\Cur,\Hom\}\), the recorded origin package
\[
\OriginPackage_{\sigma}.
\]
At benchmark-facing scope this package consists of the following data.
\begin{enumerate}
    \item An origin state space and compact convex origin body over the recorded precursor, lift, and substrate bodies, together with affine projections to those lower bodies and unique origin-to-precursor, origin-to-lift, and origin-to-substrate fiber minimizers.
    \item An origin restriction section whose zero set on the origin body is the exact origin shell.
    \item Affine observer-data and shell-response readouts factoring through the same lower readouts already used by the downstream precursor theorem.
    \item An origin shell chart to the same reservoir body, an origin shell-internal support recorder, an origin support shell, an origin support zero core, support-shell and zero-core charts, and the same charted zero/hidden equation-space splitting.
    \item A hidden charge together with the corresponding hidden recorder and the same positive hidden gap structure at benchmark scope.
    \item An observer-compatible exact origin zero-response realization.
    \item A fixed-data solvability statement and coercive control on origin data-invisible directions.
\end{enumerate}
\end{definition}

\begin{remark}
Definition \ref{def:orig} is the precise target already constructed in the recorded origin-from-ground theorem and then used by the current precursor-forcing theorem. The present note does not enlarge that target. It asks whether any benchmark-relevant freedom remains in the deeper ground layer once that exact target is required.
\end{remark}

\section{Ground Presentations and Their Benchmark-Relevant Core}

\begin{definition}[Origin-generating substrate ground dynamics]
\label{def:ground}
For each sector \(\sigma\in\{\Car,\Cur,\Hom\}\), an \emph{origin-generating substrate ground dynamics package}
\[
\GroundPackage_{\sigma}
\]
consists of the following data.
\begin{enumerate}
    \item A compact convex ground body over the recorded origin state space together with affine projections to the origin, precursor, lift, and substrate bodies and unique ground-to-origin, ground-to-precursor, ground-to-lift, and ground-to-substrate fiber minimizers.
    \item A ground restriction section whose zero set on the ground body is the exact ground shell.
    \item Affine ground observer-data and shell-response readouts factoring through the recorded origin readouts.
    \item A ground shell chart to the same reservoir body, a ground shell-internal support recorder, a ground support shell, a ground support zero core, support-shell and zero-core charts, and the same charted zero/hidden equation-space splitting.
    \item A ground hidden charge, an observer-compatible exact ground zero-response realization, fixed-data solvability, and coercive control on ground data-invisible directions.
\end{enumerate}
\end{definition}

\begin{definition}[Ground-faithful presentation]
\label{def:faithful}
A sectorwise ground presentation is \emph{ground-faithful} if its projected exact-shell transport reproduces exactly the recorded origin package of Definition \ref{def:orig}: the same origin body, the same origin-, precursor-, lift-, and substrate-fiber minimizer images, the same exact origin shell, the same shell chart to the same reservoir body, the same support shell and support zero core, the same charted zero/hidden split in the same equation variables, the same hidden charge and exact zero-response realization, the same fixed-data solvability statement, and the same coercive control. It must also leave the fixed benchmark outputs \eqref{eq:fixedoutputs} unchanged.
\end{definition}

\begin{definition}[Fiber-minimizing exact-shell ground core]
\label{def:core}
For a ground-faithful presentation, the \emph{fiber-minimizing exact-shell ground core}
\[
\GroundCore_{\sigma}
\]
is the benchmark-relevant part of the ground package consisting of:
\begin{enumerate}
    \item the ground-to-origin projection together with the ground origin-, precursor-, lift-, and substrate-fiber minimizer images;
    \item the exact ground shell and its observer-data and shell-response readouts;
    \item the ground shell chart to the same reservoir body and the ground shell-internal support recorder;
    \item the ground support shell, the ground support zero core, the support-shell and zero-core charts, and the charted zero/hidden split in the same equation variables;
    \item the hidden charge, the exact ground zero-response realization, the fixed-data solvability statement, and the coercive control on ground data-invisible directions.
\end{enumerate}
\end{definition}

\begin{definition}[Ground-faithful equivalence]
\label{def:equiv}
A \emph{ground-faithful equivalence} between two ground-faithful presentations is an identification of their cores \(\GroundCore_{\sigma}\) that preserves the same projected origin package, the same fiber-minimizer images, the same charted shell/support transport, the same zero/hidden equation-space coordinates, the same hidden quadratic form, the same exact zero-response realization, and the same solvability/coercive package.
\end{definition}

\begin{remark}
Definition \ref{def:equiv} is intentionally narrower than full off-shell identity of the ambient ground body. The active burden is not to classify every possible ground thickening outside the fiber-minimizing exact-shell core. It is to remove the benchmark-relevant freedom still visible on the active one-metric line at origin scope.
\end{remark}

\section{Necessity of the Origin Package}

\begin{proposition}[The fiber-minimizer hierarchy is necessary]
\label{prop:minimizers}
Every ground-faithful presentation determines the same origin-, precursor-, lift-, and substrate-fiber minimizer images inside the ground layer. Hence the ground-to-origin transport relevant at benchmark scope is canonical.
\end{proposition}

\begin{proof}
By Definition \ref{def:faithful}, the ground presentation must recover exactly the fixed origin package of Definition \ref{def:orig}. In particular, it must recover the same origin body and the same lower bodies together with the same unique minimizers on the corresponding recorded fibers. Since those lower fibers are fixed continuation data, the image of each ground minimizer is keyed uniquely by the lower point it lifts. The benchmark-relevant fiber-minimizer hierarchy therefore depends only on the shared lower datum and is canonical.
\end{proof}

\begin{proposition}[Exact-shell transport data are necessary]
\label{prop:shell}
Every ground-faithful presentation determines the same exact-shell transport data at benchmark scope. More precisely, the ground-to-origin projection restricts to diffeomorphic transport from the ground shell, support shell, and zero core onto the corresponding fixed origin loci, and the shell chart, support recorder, support-shell chart, zero chart, zero/hidden split, and equation readouts are thereby forced up to the canonical ground-faithful equivalence.
\end{proposition}

\begin{proof}
If a ground presentation is faithful, then by definition its projected exact-shell transport recovers precisely the shell-side data already fixed in Definition \ref{def:orig}. The ground shell must therefore project diffeomorphically to the recorded origin shell, otherwise the fixed origin shell would not be recovered as exact-shell image. The same argument applies one step further to the support shell and then to the zero core. Once those three projected loci are fixed, the shell chart to the common reservoir body, the support recorder, the support-shell chart to the equation space, the zero chart to the zero subspace, and the zero/hidden split cannot vary independently: changing any of them would change the induced origin shell transport or the fixed equation-side coordinates already recorded on the active line. Hence all these constituents are forced up to the canonical ground-faithful equivalence.
\end{proof}

\begin{proposition}[Hidden charge, zero-response realization, solvability, and coercivity are necessary]
\label{prop:hidden}
Every ground-faithful presentation determines, up to ground-faithful equivalence, the same hidden charge data, the same exact ground zero-response realization, the same fixed-data solvability statement, and the same coercive control on data-invisible directions.
\end{proposition}

\begin{proof}
The fixed origin package already records one hidden sector, one hidden quadratic form with positive gap, one exact zero-response realization, and one solvability/coercive package on the zero core. A ground-faithful presentation must project to that same origin package. Because the support-shell transport and the zero/hidden coordinates are already fixed by Proposition \ref{prop:shell}, the hidden sector cannot be changed without changing the induced origin recorder. Likewise, the exact zero-response realization must project to the same visible/current/equilibrium shell data, so its image on the ground zero core is forced by the same origin compatibility requirement. The solvability and coercive statements are properties of this same projected zero-response core; if they differed at ground scope, the induced origin statements would differ as well. Therefore those constituents are necessary up to ground-faithful equivalence.
\end{proof}

\begin{proposition}[The benchmark package remains fixed]
\label{prop:benchmark}
Every ground-faithful presentation preserves the same benchmark one-metric package \eqref{eq:fixedoutputs}. In particular, the operative metric remains exactly \(\CompletedMetric\), the ordinary matter route is unchanged, there is no second operative metric, and the low-energy matter law is not enlarged.
\end{proposition}

\begin{proof}
This is part of Definition \ref{def:faithful}. The present note removes only ground-level freedom internal to the active derivational line. It does not alter the recorded primitive outputs.
\end{proof}

\section{Main Theorem}

\begin{theorem}[Precursor-generating substrate origin dynamics necessity / uniqueness]
\label{thm:main}
Fix the primitive continuation data of Definition \ref{def:fixed} and the recorded origin package of Definition \ref{def:orig}. Then, for each sector \(\sigma\in\{\Car,\Cur,\Hom\}\), the fiber-minimizing exact-shell ground core \(\GroundCore_{\sigma}\) is necessary and unique at benchmark scope up to ground-faithful equivalence. More precisely:
\begin{enumerate}
    \item every ground-faithful presentation must determine the same origin-, precursor-, lift-, and substrate-fiber minimizer hierarchy and the same exact-shell transport data described in Definition \ref{def:core};
    \item for any two ground-faithful presentations, there is a unique ground-faithful equivalence between their benchmark-relevant cores;
    \item under that equivalence, the ground shell, support shell, zero core, shell/support/equation charts, hidden charge, exact zero-response realization, fixed-data solvability statement, and coercive control all coincide;
    \item consequently no additional benchmark-relevant origin-level freedom remains on the active one-metric line: the precursor-generating substrate origin dynamics are forced by origin-generating substrate ground dynamics at benchmark-facing scope.
\end{enumerate}
\end{theorem}

\begin{proof}
Item (1) is exactly the combined content of Propositions \ref{prop:minimizers}, \ref{prop:shell}, and \ref{prop:hidden}. Those propositions show that no ground-faithful presentation may vary the fiber-minimizing exact-shell core while still inducing the same fixed origin package.

For item (2), Proposition \ref{prop:minimizers} gives the canonical comparison on the four minimizer images, while Proposition \ref{prop:shell} gives the canonical comparison on the shell, support shell, and zero core together with their charts and recorders. Because each of those loci projects to the same fixed origin datum, the ground-faithful equivalence is unique.

For item (3), Proposition \ref{prop:shell} shows that the shell chart, support recorder, support-shell chart, zero chart, split, and equation readouts are transported identically under the canonical equivalence. Proposition \ref{prop:hidden} then shows that the hidden charge, exact zero-response realization, solvability statement, and coercive control are likewise transported identically. Hence the entire benchmark-relevant ground core agrees under the unique ground-faithful equivalence.

Item (4) is the project-level consequence of the previous items. Any remaining off-shell freedom outside the fiber-minimizing exact-shell ground core is not benchmark-relevant on the current line, because it does not change the fixed origin package or the benchmark outputs. Therefore the current honest origin burden has been removed in the only sense that matters for the active one-metric derivational program.
\end{proof}

\begin{corollary}[No remaining benchmark-relevant origin choice]
\label{cor:nofreedom}
At current scope there is no second benchmark-relevant origin presentation of the recorded origin package other than a ground-faithful relabeling of \(\GroundCore_{\sigma}\) in each sector.
\end{corollary}

\begin{proof}
If a second benchmark-relevant origin presentation existed, it would either change some constituent proved necessary in Theorem \ref{thm:main}, contradicting item (1), or else differ only by the canonical equivalence of item (2), in which case it would be merely a ground-faithful relabeling.
\end{proof}

\begin{remark}
Corollary \ref{cor:nofreedom} is the precise benchmark-facing sense in which the present manuscript is stronger than another deeper-origin relay. The current origin layer is no longer merely available because a still deeper theorem reproduces it. The benchmark-relevant ground core that yields it is now forced at current scope.
\end{remark}

\section{Routing Consequence}

\begin{corollary}[What must be done next]
\label{cor:next}
After Theorem \ref{thm:main}, the next honest primary manuscript below the current frontier must do at least one of the following:
\begin{enumerate}
    \item derive or justify the origin-generating substrate ground dynamics from still more primitive explicit substrate structure in a way that changes benchmark-facing status;
    \item identify a genuinely smaller theorem-level collapse, sharper necessity result, or sharper uniqueness result strictly inside the recorded origin package \(\OriginPackage_{\sigma}\);
    \item do either of the previous two while preserving the same benchmark one-metric package, the same ordinary matter route, and the same claim boundary.
\end{enumerate}
It is no longer enough merely to restate the current origin package, the current precursor package \(\PrecursorPackage\), the current lift package \(\LiftPackage\), or the already forced shell-restriction package \(\ShellRestriction\), and it is no longer enough merely to insert another same-output relay above the fixed ground core.
\end{corollary}

\begin{proof}
Theorem \ref{thm:main} removes the benchmark-relevant freedom presently visible at origin level. The next honest burden therefore lies either below that level or in a strictly smaller internal compression of it. Another package-preserving restatement of the same origin package would not change project status.
\end{proof}

\section{Conclusion}

The positive result of this note is not that a final substrate microphysics has already been found. Its gain is narrower and more honest. The recorded origin-from-ground theorem had already shown that the current origin package is the canonical projected exact-shell output of a deeper ground class. The present theorem then removes the remaining benchmark-relevant freedom inside that origin layer itself. At current scope, the fiber-minimizing exact-shell ground core is necessary and unique up to canonical ground-faithful equivalence.

That conclusion is exactly the sharper theorem-level gain the project needed. The ground-to-origin projection, the four fiber-minimizer images, the ground shell and support transport, the zero/hidden split in the same equation variables, the hidden charge, the exact zero-response realization, and the solvability/coercive package are no longer merely one possible deeper presentation of the same downstream line. Once the presentation is required to induce the fixed origin package, those constituents are forced.

The benchmark boundary remains unchanged. The same primitive shell remains the shell carrying the observer outputs, so the one operative metric is still exactly \(\CompletedMetric\), the ordinary matter route is unchanged, there is no second operative metric, and the low-energy matter law is not enlarged. Nothing here upgrades adoption to derivation or licenses stronger first-principles promotion language.

The honest open burden now lies below origin scope at ground scope. The next benchmark-facing manuscript should therefore first try to derive or uniquely force the origin-generating substrate ground dynamics from still more primitive explicit substrate structure in a way that actually changes benchmark-facing status. If that cannot yet be done cleanly, the admissible fallback is a genuinely smaller collapse, sharper necessity result, or sharper uniqueness result strictly inside the precursor-generating substrate origin package while preserving the same benchmark one-metric package and claim boundary.

\end{document}
